\documentclass[11pt]{article}
\usepackage{amsmath,amssymb,amsthm}
\usepackage[margin=1in]{geometry}
\usepackage[T1]{fontenc}
\usepackage{hyperref}

\newtheorem{theorem}{Theorem}[section]
\newtheorem{lemma}[theorem]{Lemma}
\newtheorem{proposition}[theorem]{Proposition}
\newtheorem{corollary}[theorem]{Corollary}
\newtheorem{definition}[theorem]{Definition}
\theoremstyle{remark}
\newtheorem{remark}[theorem]{Remark}

\DeclareMathOperator{\Cas}{W}

\title{Solution to Ramanujan Challenge Problem 2.3:\\
The Sum $\pi + e$ as an Ap\'ery Limit}
\author{Xiang Huang\\{\small University of Illinois Springfield}\\{\small\texttt{xhuan5@uis.edu}}}
\date{August 2026}

\begin{document}
\maketitle

\begin{abstract}
We prove that the order-four recurrence of Problem~2.3 converges to $\pi+e$.
The mechanism is that the challenge operator is a \emph{tensor product}: order
$4 = 2\times 2$. It annihilates every product $X_{n+2}Y_{n+3}$ of a solution
$X$ of Lambert's continuant recurrence (whose ratio produces $\pi/4$) with a
solution $Y$ of the derangement recurrence (whose ratio produces $e$). The key
observation is that $m!$ satisfies the \emph{same} recurrence as the derangement
numbers $D_m$; this is why $\pi$ and $e$ appear together, and additively. The
challenge's own $p_n,q_n$ are then identified with explicit tensor solutions,
and the ratio splits \emph{exactly}, with no asymptotic analysis:
\[
\frac{p_n}{q_n} \;=\; 4\,\frac{B_{n+2}}{A_{n+2}} \;+\; \frac{(n+3)!}{D_{n+3}}.
\]
Every step is proved here and machine-checked in Lean~4. In particular, the
Lambert limit $B_m/A_m\to\pi/4$ is derived from a positive moment integral and
a geometric remainder bound; it is not cited as an external input.
\end{abstract}

\section{The problem}

For $n \ge 1$ the challenge poses the recurrence
\begin{equation}\label{eq:rec}
\begin{aligned}
0 = {}& \bigl(-n^3 + 2n^2 + 7n + 3\bigr)\,u_n \\
 &+ (n+2)\bigl(2n^4 + n^3 - 26n^2 - 48n - 19\bigr)\,u_{n-1} \\
 &+ (n+2)\bigl(n^6 + 9n^5 + 8n^4 - 87n^3 - 249n^2 - 234n - 68\bigr)\,u_{n-2} \\
 &+ (n+1)^2(n+2)\bigl(2n^5 + 3n^4 - 13n^3 - 21n^2 + 4\bigr)\,u_{n-3} \\
 &- n^3(n+1)^2(n+2)\bigl(n^3 + n^2 - 8n - 11\bigr)\,u_{n-4},
\end{aligned}
\end{equation}
and asks for $\lim_{n\to\infty} p_n/q_n = \pi+e$, where $p,q$ are the solutions
with
\begin{equation}\label{eq:init}
(p_{-3},p_{-2},p_{-1},p_0) = (1,1,20,296),
\qquad
(q_{-3},q_{-2},q_{-1},q_0) = (1,0,4,48).
\end{equation}
Write $c_0(n),\dots,c_4(n)$ for the five coefficient polynomials
in~\eqref{eq:rec}, in the order displayed, so that~\eqref{eq:rec} reads
$\sum_{j=0}^{4} c_j(n)\,u_{n-j}=0$.

\section{Two order-two systems}

\begin{definition}
The \emph{Lambert system} $\mathcal{L}$ is the recurrence
\begin{equation}\label{eq:lambert}
X_m = (2m+1)X_{m-1} + m^2 X_{m-2} \qquad (m \ge 1),
\end{equation}
and the \emph{derangement system} $\mathcal{D}$ is
\begin{equation}\label{eq:derange}
Y_m = (m-1)\bigl(Y_{m-1} + Y_{m-2}\bigr) \qquad (m \ge 2).
\end{equation}
Distinguished solutions: $A_m,B_m$ of $\mathcal{L}$ with
$A_{-1}=A_0=1$ and $B_{-1}=0$, $B_0=1$; and $D_m$ (the derangement numbers) of
$\mathcal{D}$ with $D_0=1,\ D_1=0$.
\end{definition}

Thus $A_1=4$, $A_2=24$, $A_3=204$; $B_1=3$, $B_2=19$, $B_3=160$; and
$D_2=1$, $D_3=2$, $D_4=9$.

\begin{lemma}[The factorial is a second solution of $\mathcal{D}$]\label{lem:fact}
$F_m := m!$ satisfies~\eqref{eq:derange}.
\end{lemma}

\begin{proof}
$(m-1)\bigl((m-1)! + (m-2)!\bigr) = (m-2)!\,(m-1)\bigl((m-1)+1\bigr)
= (m-2)!\,(m-1)m = m!$.
\end{proof}

Lemma~\ref{lem:fact} is the structural reason the constant $e$ occurs in this
problem, and the reason it occurs \emph{additively} alongside $\pi$: the
$\mathcal{D}$-system carries both $D_m$ and $m!$, and their ratio is the
classical approximation to $e$.

\begin{lemma}[Casorati determinants]\label{lem:cas}
For all $m\ge 0$,
\[
A_mB_{m-1}-A_{m-1}B_m = (-1)^{m+1}(m!)^2,
\qquad
D_mF_{m-1}-D_{m-1}F_m = (-1)^{m}(m-1)!\ \ (m\ge1).
\]
In particular $\{A,B\}$ is a basis of the solution space of $\mathcal{L}$ and
$\{D,F\}$ is a basis of the solution space of $\mathcal{D}$.
\end{lemma}

\begin{proof}
Write $\Cas_m = A_mB_{m-1}-A_{m-1}B_m$. Substituting~\eqref{eq:lambert} for
$A_m$ and $B_m$, the $(2m+1)$-terms cancel and
$\Cas_m = -m^2\,\Cas_{m-1}$; with $\Cas_0 = A_0B_{-1}-A_{-1}B_0 = -1$ this
telescopes to $\Cas_m = (-1)^{m+1}(m!)^2$. The second identity is the same
computation: substituting~\eqref{eq:derange} gives
$D_mF_{m-1}-D_{m-1}F_m = -(m-1)\bigl(D_{m-1}F_{m-2}-D_{m-2}F_{m-1}\bigr)$, and
the initial value at $m=1$ is $D_1F_0-D_0F_1 = -1$.
\end{proof}

\section{The tensor theorem}

\begin{theorem}[Tensor annihilation]\label{thm:tensor}
Let $X$ be \emph{any} solution of $\mathcal{L}$ and $Y$ \emph{any} solution of
$\mathcal{D}$. Then
\[
u_n := X_{n+2}\,Y_{n+3}
\]
satisfies the challenge recurrence~\eqref{eq:rec} for every $n\ge1$.
\end{theorem}

\begin{proof}
The assertion is an identity in the four free initial values
$X_{-1},X_0,Y_0,Y_1$, so it suffices to prove a polynomial identity. Put
\[
\mathfrak{a}(m) = \begin{pmatrix} 2m+1 & m^2\\ 1 & 0\end{pmatrix},
\qquad
\mathfrak{d}(m) = \begin{pmatrix} m-1 & m-1\\ 1 & 0\end{pmatrix},
\]
so that $(X_m,X_{m-1})^{\mathsf T} = \mathfrak{a}(m)(X_{m-1},X_{m-2})^{\mathsf T}$
and $(Y_m,Y_{m-1})^{\mathsf T} = \mathfrak{d}(m)(Y_{m-1},Y_{m-2})^{\mathsf T}$.
Then the vector
\[
v_n := (X_{n+2},X_{n+1})^{\mathsf T}\otimes(Y_{n+3},Y_{n+2})^{\mathsf T}\in\mathbb{Q}^4
\]
obeys $v_n = T(n)\,v_{n-1}$ with $T(n) = \mathfrak{a}(n+2)\otimes\mathfrak{d}(n+3)$,
and $u_n = e_1^{\mathsf T}v_n$ where $e_1$ is the first standard covector. Hence
\[
\sum_{j=0}^{4} c_j(n)\,u_{n-j}
= \Biggl(\sum_{j=0}^{4} c_j(n)\,e_1^{\mathsf T}\,T(n)T(n-1)\cdots T(n-j+1)\Biggr) v_{n-4},
\]
the empty product being the identity. As the initial data vary, $v_{n-4}$ ranges
over a spanning set of $\mathbb{Q}^4$ (Lemma~\ref{lem:cas} supplies four
independent products), so the recurrence holds for all solutions if and only if
the covector in parentheses vanishes identically --- that is, if and only if
four polynomial identities in $\mathbb{Z}[n]$ hold.

Concretely, one eliminates $X_{n+2},X_{n+1},X_n$ and $Y_{n+3},Y_{n+2},Y_{n+1}$
using~\eqref{eq:lambert} and~\eqref{eq:derange} until $\sum_j c_j(n)u_{n-j}$ is
a bilinear form in the four products
\[
X_{n-2}Y_{n-1},\quad X_{n-2}Y_{n},\quad X_{n-1}Y_{n-1},\quad X_{n-1}Y_{n},
\]
whose four coefficients are polynomials in $n$. All four are identically zero.

This is a finite computation in $\mathbb{Z}[n]$, and it has been carried out
three independent ways: (i) in Lean~4, where after substitution it is
discharged by the \texttt{ring} normaliser and is therefore kernel-checked
(\texttt{RamanujanChallenge.P23.tensor\_rec}); (ii) symbolically over
$\mathbb{Q}[n]$ with four free initial values; and (iii) in exact integer
arithmetic for the concrete sequences. See~\S\ref{sec:formal}
and~\S\ref{sec:numeric}.
\end{proof}

\begin{corollary}\label{cor:span}
The solution space of~\eqref{eq:rec} over $\mathbb{Q}$ is exactly four
dimensional and is spanned by
\[
\{A_{n+2}D_{n+3},\ A_{n+2}F_{n+3},\ B_{n+2}D_{n+3},\ B_{n+2}F_{n+3}\}.
\]
Thus the challenge operator \emph{is} the tensor product of $\mathcal{L}$
and $\mathcal{D}$, not merely annihilated by it.
\end{corollary}

\begin{proof}
By Lemma~\ref{lem:cas} the four listed sequences are linearly independent: their
value matrix at two consecutive indices is the Kronecker product of two
invertible $2\times2$ matrices. By Theorem~\ref{thm:tensor} they solve
\eqref{eq:rec}. Since $c_0(n)\neq0$ for $n\ge1$ (Lemma~\ref{lem:c0} below), the
solution space of~\eqref{eq:rec} has dimension exactly four.
\end{proof}

\section{Identification of the challenge sequences}

\begin{lemma}\label{lem:c0}
$c_0(n) = -n^3+2n^2+7n+3 \ne 0$ for every integer $n\ge1$; consequently a
solution of~\eqref{eq:rec} is uniquely determined by
$(u_{-3},u_{-2},u_{-1},u_0)$.
\end{lemma}

\begin{proof}
For $n=1,2,3,4$ the values are $11,17,15,-1$. For $n\ge5$,
$n^3 - 2n^2 - 7n - 3 = n^2(n-2) - 7n-3 \ge 3n^2 - 7n - 3 > 0$, so
$c_0(n)<0$. As $c_0(n)\ne0$, equation~\eqref{eq:rec} solves for $u_n$ at each
$n\ge1$.
\end{proof}

\begin{proposition}[Closed forms]\label{prop:closed}
The challenge sequences of~\eqref{eq:init} are
\begin{equation}\label{eq:closed}
q_n = A_{n+2}\,D_{n+3},
\qquad
p_n = 4\,B_{n+2}\,D_{n+3} + A_{n+2}\,(n+3)!\, .
\end{equation}
\end{proposition}

\begin{proof}
By Theorem~\ref{thm:tensor} the sequence $A_{n+2}D_{n+3}$ solves~\eqref{eq:rec};
so do $B_{n+2}D_{n+3}$ and, by Lemma~\ref{lem:fact}, $A_{n+2}(n+3)!$. Since
\eqref{eq:rec} is linear, the right-hand sides of~\eqref{eq:closed} are
solutions. It remains to check four initial values of each, using
$A_{-1}=A_0=1,A_1=4,A_2=24$, $B_{-1}=0,B_0=1,B_1=3,B_2=19$ and
$D_0=1,D_1=0,D_2=1,D_3=2$:
\[
\begin{array}{llll}
A_{-1}D_0 = 1, & A_0D_1 = 0, & A_1D_2 = 4, & A_2D_3 = 48;\\[2pt]
4B_{-1}D_0 + A_{-1}0! = 1, &
4B_0D_1 + A_0 1! = 1, &
4B_1D_2 + A_1 2! = 20, &
4B_2D_3 + A_2 3! = 296.
\end{array}
\]
These are exactly~\eqref{eq:init}. By Lemma~\ref{lem:c0} the solution with
prescribed initial data is unique, so~\eqref{eq:closed} holds for all $n$.
\end{proof}

\begin{corollary}[Exact splitting]\label{cor:split}
For every $n \ge -1$,
\begin{equation}\label{eq:split}
\frac{p_n}{q_n} \;=\; 4\,\frac{B_{n+2}}{A_{n+2}} \;+\; \frac{(n+3)!}{D_{n+3}}.
\end{equation}
\end{corollary}

\begin{proof}
Divide~\eqref{eq:closed}. The denominators are nonzero: $A_m>0$ for all $m$ by
induction on~\eqref{eq:lambert} (both coefficients are positive for $m\ge1$),
and $D_m>0$ for $m\ge2$ likewise.
\end{proof}

We stress that~\eqref{eq:split} is an identity, valid term by term; no
asymptotic or Poincar\'e--Perron argument is used to obtain it. All that
remains is to evaluate the two limits separately.

\section{The two limits}

\begin{lemma}[The $e$ half]\label{lem:e}
$\displaystyle \lim_{m\to\infty}\frac{m!}{D_m} = e$.
\end{lemma}

\begin{proof}
Write $S_m := \sum_{k=0}^{m}(-1)^k/k!$. We claim $D_m = m!\,S_m$ for all
$m\ge0$. This holds at $m=0$ ($1 = 1\cdot1$) and $m=1$ ($0 = 1\cdot(1-1)$).
Assume it at $m-1$ and $m-2$. Using $(m-1)! = (m-1)(m-2)!$,
\[
(m-1)\bigl(D_{m-1}+D_{m-2}\bigr)
= (m-1)\bigl[(m-1)!\,S_{m-1} + (m-2)!\,S_{m-2}\bigr]
= (m-1)!\bigl[(m-1)S_{m-1} + S_{m-2}\bigr].
\]
Now $S_{m-2} = S_{m-1} - (-1)^{m-1}/(m-1)!$, so the bracket equals
$m\,S_{m-1} - (-1)^{m-1}/(m-1)!$ and therefore
\[
(m-1)\bigl(D_{m-1}+D_{m-2}\bigr)
= m!\,S_{m-1} - (-1)^{m-1}
= m!\,S_{m-1} + (-1)^{m}
= m!\Bigl(S_{m-1} + \frac{(-1)^m}{m!}\Bigr)
= m!\,S_m,
\]
which is $D_m$. Hence $m!/D_m = S_m^{-1}$. The partial sums $S_m$ converge to
$e^{-1}\neq0$, and $S_m>0$ for $m\ge2$ because an alternating series with
strictly decreasing terms has all its partial sums from index~$2$ on trapped
between $S_2 = 1/2$ and $S_3 = 1/3$. So the reciprocals converge to $e$.
\end{proof}

\begin{theorem}[The $\pi$ half; Lambert]\label{thm:lambert}
$\displaystyle \lim_{m\to\infty}\frac{B_m}{A_m} = \frac{\pi}{4}$.
\end{theorem}

\begin{proof}
Put
\[
s=\sqrt2,\qquad c=2-s,\qquad P(x)=x(1-x),\qquad D(x)=1-cP(x),
\]
and, for $r\ge0$, define the positive moments
\begin{equation}\label{eq:moment}
K_r:=\int_0^1\frac{P(x)^r}{D(x)^{r+1}}\,dx.
\end{equation}
On $[0,1]$ we have $0\le P(x)\le1/4$ and $0<c<1$, hence $D(x)>0$.

We first compute two initial moments. Substituting $y=2x-1$ and using
$\arctan(s-1)=\pi/8$ gives
\begin{equation}\label{eq:K0}
K_0=\frac4s\arctan(s-1)=s\frac{\pi}{4}.
\end{equation}
For completeness, the arctangent identity follows from the double-angle
formula: $2(s-1)/(1-(s-1)^2)=1$, and both angles lie in $(0,\pi/2)$.
Next, direct differentiation gives
\[
\frac{d}{dx}\frac{1-2x}{D(x)}
=-s\frac1{D(x)}-2\frac{P(x)}{D(x)^2}.
\]
The expression on the left takes the values $1$ and $-1$ at the two
endpoints. Integrating and using~\eqref{eq:K0}, we obtain
\begin{equation}\label{eq:K1}
K_1=1-\frac{\pi}{4}.
\end{equation}

The same integration-by-parts mechanism gives the whole recurrence. For every
$r\ge0$, algebraic differentiation yields
\begin{align*}
\frac{d}{dx}\left(
  \frac{(1-2x)P(x)^{r+1}}{D(x)^{r+2}}
\right)
={}&(r+1)\frac{P(x)^r}{D(x)^{r+1}}
 -(2r+3)s\frac{P(x)^{r+1}}{D(x)^{r+2}}\\
&-2(r+2)\frac{P(x)^{r+2}}{D(x)^{r+3}}.
\end{align*}
Because $P(0)=P(1)=0$, the boundary term vanishes. Thus
\begin{equation}\label{eq:Krec}
2(r+2)K_{r+2}=(r+1)K_r-(2r+3)sK_{r+1}.
\end{equation}

Define
\[
E_r:=\frac{r!s^r}{s}K_r.
\]
Equations~\eqref{eq:K0}--\eqref{eq:Krec} become
\[
E_0=\frac{\pi}{4},\qquad E_1=1-\frac{\pi}{4},\qquad
E_{r+2}=(r+1)^2E_r-(2r+3)E_{r+1}.
\]
Comparing this recurrence and its two initial values with
\eqref{eq:lambert}, induction gives the exact remainder formula
\begin{equation}\label{eq:lambert-remainder}
A_{r-1}\frac{\pi}{4}-B_{r-1}=(-1)^r E_r
\qquad(r\ge0).
\end{equation}

It remains only to show that this remainder is negligible relative to $A_m$.
Since $P/D\le1/(2+s)$ on $[0,1]$, positivity of the integrand gives
\[
0\le K_r\le \left(\frac1{2+s}\right)^rK_0.
\]
Using $s/(2+s)=s-1$ and~\eqref{eq:K0}, this implies
\begin{equation}\label{eq:Ebound}
0\le E_r\le r!(s-1)^r\frac{\pi}{4}.
\end{equation}
The positive recurrence for $A_m$ also gives $A_{r-1}\ge r!$ by induction.
Taking $r=m+1$ in~\eqref{eq:lambert-remainder} and applying
\eqref{eq:Ebound}, we conclude that
\[
\left|\frac{\pi}{4}-\frac{B_m}{A_m}\right|
\le \frac{\pi}{4}(\sqrt2-1)^{m+1}\longrightarrow0,
\]
because $0<\sqrt2-1<1$.
\end{proof}

\begin{theorem}[Main]\label{thm:main}
$\displaystyle\lim_{n\to\infty}\frac{p_n}{q_n} = \pi + e$.
\end{theorem}

\begin{proof}
Combine the exact splitting~\eqref{eq:split} with Lemma~\ref{lem:e} and
Theorem~\ref{thm:lambert}:
\[
\frac{p_n}{q_n} = 4\,\frac{B_{n+2}}{A_{n+2}} + \frac{(n+3)!}{D_{n+3}}
\;\longrightarrow\; 4\cdot\frac{\pi}{4} + e = \pi + e. \qedhere
\]
\end{proof}

\begin{remark}[Why the problem looks hard]
The four Poincar\'e behaviours of~\eqref{eq:rec} are the products of the two
$\mathcal{L}$-behaviours (after a factorial gauge, $1\pm\sqrt2$) with the two
$\mathcal{D}$-behaviours. The dominant solution is not the one carrying the
answer; $p_n$ and $q_n$ are particular mixtures whose ratio has the two
constants sitting in \emph{different} tensor factors. Any attempt to read
$\pi+e$ off the order-four recurrence directly must therefore separate a
rank-two continuant factor from a rank-two factorial factor, which is what
Theorem~\ref{thm:tensor} does in closed form.
\end{remark}

\section{Formal verification}\label{sec:formal}

The accompanying Lean~4 development (\texttt{lean/}, toolchain
\texttt{leanprover/lean4:v4.29.0}, Mathlib pinned in \texttt{lake-manifest.json})
builds with \textbf{0 errors and 0 \texttt{sorry}} and formalises the entire
argument above, including Theorem~\ref{thm:lambert}. The principal statements are:

\begin{center}
\begin{tabular}{ll}
\texttt{tensor\_rec} & Theorem~\ref{thm:tensor}, for arbitrary $X,Y$\\
\texttt{factorial\_isDerRec} & Lemma~\ref{lem:fact}\\
\texttt{challengeQ\_rec}, \texttt{challengeP\_rec} & the closed forms solve~\eqref{eq:rec}\\
\texttt{challengeQ\_zero}\,\dots\,\texttt{challengeP\_three} & the eight initial values~\eqref{eq:init}\\
\texttt{C0\_ne\_zero}, \texttt{eq\_of\_satisfiesRec} & Lemma~\ref{lem:c0} and uniqueness\\
\texttt{ratio\_split} & Corollary~\ref{cor:split}\\
\texttt{lambertB\_div\_lambertA\_tendsto\_pi\_div\_four} & Theorem~\ref{thm:lambert}\\
\texttt{factorial\_div\_derang\_tendsto\_exp\_one} & Lemma~\ref{lem:e}\\
\texttt{problem23\_pi\_add\_e} & Theorem~\ref{thm:main}\\
\end{tabular}
\end{center}

Every one of these depends only on the three standard axioms
\texttt{propext}, \texttt{Classical.choice}, \texttt{Quot.sound}; the purely
algebraic ones (\texttt{tensor\_rec}, \texttt{challengeQ\_rec},
\texttt{challengeP\_rec}) do not even use \texttt{Classical.choice}. There is no
\texttt{native\_decide} anywhere. The $e$ half is obtained from Mathlib's
\texttt{numDerangements\_tendsto\_inv\_e} after proving that our \texttt{derang}
agrees with Mathlib's \texttt{numDerangements}.

The top-level theorem has no mathematical hypotheses:
\begin{verbatim}
theorem problem23_pi_add_e :
    Tendsto (fun m => (challengeP m : R) / (challengeQ m : R))
      atTop (nhds (Real.pi + Real.exp 1))
\end{verbatim}

\section{Numerical verification}\label{sec:numeric}

The script \texttt{verify.py} performs three checks in exact arithmetic.
\begin{enumerate}
\item It solves~\eqref{eq:rec} forward over $\mathbb{Q}$ from the initial
data~\eqref{eq:init} alone and confirms that the result is integral and agrees
with the closed forms~\eqref{eq:closed}. The first new values are
$q_1 = 1836$, $q_2 = 97680$, $p_1 = 10656$, $p_2 = 573344$.
\item It verifies Theorem~\ref{thm:tensor} symbolically, with four free initial
values, so that no accidental agreement of a particular solution can be
mistaken for the general identity.
\item It evaluates $p_n/q_n$ in high precision: at $n=50$ the agreement with
$\pi+e = 5.859874482048838\ldots$ is to $40$ decimal places.
\end{enumerate}

\begin{thebibliography}{1}
\bibitem{Challenge2026}
Ramanujan Machine Group,
\emph{The Ramanujan Challenge for AI}, 2026.

\end{thebibliography}

\end{document}
